\documentclass[11pt]{article}
\usepackage[utf8]{inputenc}
\usepackage[T1]{fontenc}
\usepackage{amsmath,amssymb,amsfonts}
\usepackage{booktabs}
\usepackage{graphicx}
\usepackage{hyperref}
\usepackage{multirow}
\usepackage{xcolor}
\usepackage[margin=1in]{geometry}
\usepackage{natbib}
\usepackage{caption}

\title{What Controls the Discriminability of Recursive Synthetic Text?\\An Empirical Analysis via Data Processing Inequality}

\author{
Anonymous Authors
}

\date{}

\begin{document}
\maketitle

% ============================================================
\begin{abstract}
As large language models are increasingly deployed to generate web content, corpora used for pretraining inevitably contain multi-generational recursive synthetic text---outputs generated by models trained or prompted on previous synthetic outputs. We formalize recursive generation as an iterated lossy channel and invoke the Data Processing Inequality (DPI) to prove that a finite discrimination boundary $K^*$ must exist: beyond generation depth $K^*$, no statistical method can reliably distinguish consecutive depths. We then conduct the first systematic empirical measurement of $K^*$ across controlled configurations. Contrary to prior assumptions of a stable universal boundary, we find that $K^*$ is \emph{configuration-dependent}: it varies from 1 to 3 as a function of decoding strategy, model scale, recursion mechanism, and feature representation. Greedy decoding yields $K^*{=}1$; high-temperature sampling ($T{=}1.2$) extends it to $K^*{=}3$; scaling from 1.5B to 7B parameters reduces $K^*$ from 2 to 1. Learned representations push $K^*$ one level beyond hand-crafted features, revealing that measured $K^*$ partially reflects feature engineering limits rather than signal absence. We further show that graduated depth-aware filtering does not significantly outperform binary filtering ($\Delta{<}0.5$pp, $p{>}0.10$), consistent with the limited ordinal resolution at low $K^*$. Our results reframe $K^*$ from a stable property of recursive text to an operational parameter controlled by generation configuration, with direct implications for synthetic data governance.
\end{abstract}

% ============================================================
\section{Introduction}

The recursive contamination of web corpora---where language models generate text that is later scraped, mixed into training data, and used to train successor models---has emerged as a central concern for data governance \citep{shumailov2024collapse,alemohammad2024selfconsuming}. The resulting ``model collapse'' phenomenon, in which distributional diversity erodes across generations, has been studied at the macroscopic level. Yet a fundamental question remains underexplored: \emph{at what generational depth does recursive synthetic text become statistically indistinguishable from its predecessor?}

Current detection methods treat synthetic text as a binary category: real or generated. Detectors such as DivEye and Binoculars achieve high accuracy on 0-vs-1 discrimination but do not estimate \emph{how many rounds} of recursive generation a document has undergone. Data selection methods (DSIR, D4, DoReMi) perform distribution matching without depth awareness. Drayson et al.\ apply binary detection scores for importance resampling to mitigate model collapse, the closest prior work to depth-aware curation.

We formalize recursive generation as an iterated lossy channel and invoke the Data Processing Inequality (DPI) to establish that a finite discrimination boundary $K^*$ must exist: the mutual information between consecutive generation depths monotonically decreases, guaranteeing convergence to indistinguishability. We then measure $K^*$ empirically under controlled conditions.

Our central finding is negative with respect to universality: \textbf{$K^*$ is not a stable property of recursive text, but a configuration-dependent operational parameter.} We map the factors that control $K^*$:
\begin{itemize}
    \item \textbf{Decoding strategy}: $K^*$ ranges from 1 (greedy) to 3 ($T{=}1.2$) for the same model (\S\ref{sec:decoding}).
    \item \textbf{Model scale}: $K^*$ drops from 2 (1.5B) to 1 (7B) as model capacity increases (\S\ref{sec:scale}).
    \item \textbf{Recursion mechanism}: Retrain-chains converge faster than prompt-chains, reducing $K^*$ (\S\ref{sec:retrain}).
    \item \textbf{Feature representation}: Learned embeddings extend $K^*$ by one level beyond hand-crafted features (\S\ref{sec:learned}).
\end{itemize}

We further report that graduated depth-aware filtering---which assigns weights based on estimated generation depth rather than a binary real/synthetic decision---does not significantly outperform binary filtering under the conditions tested (\S\ref{sec:filtering}). This negative result is consistent with the limited ordinal resolution when $K^*$ is low.

Our contributions are: (1)~a DPI-based theoretical framework establishing the necessity of finite $K^*$; (2)~the first systematic empirical characterization of factors controlling $K^*$; (3)~evidence that measured $K^*$ partially reflects feature engineering limits; and (4)~a negative result on graduated filtering that constrains the practical utility of depth-aware curation at current resolution levels.

% ============================================================
\section{Background and Theoretical Framework}
\label{sec:background}

\subsection{Recursive Generation as Iterated Lossy Channel}

Let $X_0$ denote human-written text drawn from distribution $P_0$. A language model $\mathcal{M}$ defines a stochastic channel $T: \mathcal{X} \to \mathcal{X}$ mapping input text to generated output. Recursive generation produces a Markov chain:
\[
X_0 \xrightarrow{T} X_1 \xrightarrow{T} X_2 \xrightarrow{T} \cdots
\]
where $X_k$ denotes generation-$k$ text. Because $\mathcal{M}$ is an imperfect channel (it cannot perfectly reproduce its input distribution), each application of $T$ constitutes a lossy transformation.

\subsection{Data Processing Inequality and $K^*$}

By the Data Processing Inequality, for any Markov chain $X_0 \to X_1 \to X_2$:
\[
I(X_0; X_2) \leq I(X_0; X_1)
\]
More generally, the mutual information $I(X_k; X_{k+1})$ is non-increasing in $k$. This guarantees that discriminability between consecutive generations monotonically decreases. We define the \emph{empirical discrimination boundary} $K^*$ as the largest $k$ such that a classifier can reliably distinguish generation $k$ from generation $k{+}1$:
\[
K^* = \max\{k : \text{AUC}(k \text{ vs } k{+}1) > \tau\}
\]
where $\tau$ is a reliability threshold (we use $\tau = 0.60$ with bootstrap 95\% CI lower bound exceeding $\tau$). DPI guarantees $K^* < \infty$; we measure its value empirically.

The measured $K^*$ depends on the feature representation used for classification. We denote $K^*_f$ the boundary measured with feature set $f$, acknowledging that $K^*_f \leq K^*_{\text{optimal}}$ where the optimal uses all available information.

\subsection{Related Work}

\paragraph{Model collapse.} Shumailov et al.\ and Alemohammad et al.\ demonstrate that iterative training on self-generated data leads to distributional collapse. Dohmatob et al.\ provide theoretical analysis of convergence rates. These works characterize the \emph{phenomenon} but do not provide \emph{detection tools} for generation depth.

\paragraph{Binary synthetic text detection.} DivEye, Binoculars, DetectGPT, and Min-K\%++ achieve high accuracy on real-vs-synthetic classification. However, their confidence scores correlate only weakly with generation depth (we verify this in \S\ref{sec:binary_baseline}).

\paragraph{Data selection.} DSIR performs importance resampling to match a target distribution; Drayson et al.\ extend binary detection for data curation. Neither estimates generation depth.

% ============================================================
\section{Method}
\label{sec:method}

\subsection{Depth Chain Construction}

We construct recursive depth chains using two mechanisms:

\paragraph{Prompt-chain (canonical).} Given a source document $x_0$, we prompt model $\mathcal{M}$ to generate a continuation $x_1$. Then $x_1$ is used as context to generate $x_2$, and so on up to depth 4. We use nucleus sampling ($p=0.95$) as the default decoding strategy unless otherwise specified.

\paragraph{Retrain-chain.} We fine-tune model $\mathcal{M}$ on generation-$k$ text using LoRA (rank 16), then generate new text from the fine-tuned model to produce generation $k{+}1$. This more closely mimics real-world model collapse dynamics.

\subsection{Feature Extraction}

We extract a 15-dimensional hand-crafted feature vector from each document, comprising:
\begin{itemize}
    \item \textbf{Surprisal statistics} (6 features): mean, standard deviation, and first/second-order difference statistics of token-level surprisal under a reference model.
    \item \textbf{Lexical statistics} (5 features): vocabulary richness (type-token ratio), frequency entropy, hapax legomena ratio, and related measures.
    \item \textbf{Structural statistics} (4 features): sentence length variance, punctuation density, repetition metrics, and tail-distribution measures.
\end{itemize}

As a comparison, we also evaluate \textbf{learned representations}: frozen embeddings from Qwen2.5-1.5B followed by a 2-layer MLP classifier.

\subsection{$K^*$ Measurement Protocol}

For each pair of consecutive depths $(k, k{+}1)$, we train a binary classifier (LightGBM) on the feature vectors and evaluate via 5-fold cross-validation AUC. We compute bootstrap 95\% confidence intervals (10,000 resamples) and declare $K^* = k$ if $\text{AUC}_{k \text{vs} k{+}1}$ has CI lower bound $> 0.60$ and $\text{AUC}_{k{+}1 \text{vs} k{+}2}$ has CI upper bound $< 0.60$.

% ============================================================
\section{Experiments}
\label{sec:experiments}

Unless otherwise noted, experiments use C4 documents, Qwen2.5-1.5B in continuation mode, nucleus sampling $p{=}0.95$, 5K source documents, and hand-crafted 15D features.

% ------------------------------------------------------------
\subsection{$K^*$ Depends on Decoding Strategy}
\label{sec:decoding}

We generate depth chains under four decoding strategies---greedy, nucleus $p{=}0.9$, temperature $T{=}0.7$, and temperature $T{=}1.2$---in addition to the baseline nucleus $p{=}0.95$.

\begin{table}[h]
\centering
\caption{$K^*$ as a function of decoding strategy (Qwen2.5-1.5B, continuation mode). Higher sampling entropy yields larger $K^*$.}
\label{tab:decoding}
\begin{tabular}{lccccc}
\toprule
Decoding & $K^*$ & AUC 1v2 & 95\% CI & AUC 2v3 & AUC 3v4 \\
\midrule
Greedy & 1 & 0.572 & [0.563, 0.581] & 0.538 & 0.512 \\
T=0.7 & 1--2 & 0.591 & [0.581, 0.601] & 0.546 & 0.519 \\
Nucleus $p$=0.95 (baseline) & 2 & 0.611 & [0.602, 0.621] & 0.561 & 0.541 \\
Nucleus $p$=0.9 & 2 & 0.628 & [0.618, 0.637] & 0.554 & 0.517 \\
T=1.2 & 3 & 0.683 & [0.673, 0.694] & 0.612 & 0.558 \\
\bottomrule
\end{tabular}
\end{table}

$K^*$ varies monotonically with sampling entropy: greedy decoding yields $K^*{=}1$ (1v2 AUC CI upper $= 0.581 < 0.60$), while $T{=}1.2$ extends $K^*$ to 3 (2v3 AUC $= 0.612$, CI lower $= 0.602 > 0.60$). The temperature $T{=}0.7$ condition is ambiguous, with the 1v2 CI straddling 0.60. This demonstrates that $K^*$ is not a fixed property of the model but an operational parameter controlled by decoding configuration. Feature ranking consistency across strategies (Spearman $\rho = 0.78$) indicates that the same degradation mechanisms operate at different rates.

% ------------------------------------------------------------
\subsection{$K^*$ Depends on Model Scale}
\label{sec:scale}

We compare Qwen2.5-1.5B and Qwen2.5-7B in continuation mode on 1K C4 documents, generating depth chains to depth 3.

\begin{table}[h]
\centering
\caption{Scale effect on $K^*$. Larger models produce recursive text that converges faster to the stationary distribution.}
\label{tab:scale}
\begin{tabular}{lccccr}
\toprule
Model & Scale & $K^*$ & AUC 1v2 & 95\% CI & AUC 2v3 \\
\midrule
Qwen2.5-1.5B (cont) & 1.5B & 2 & 0.611 & [0.602, 0.621] & 0.561 \\
Qwen2.5-7B (cont) & 7B & 1 & 0.577 & [0.558, 0.596] & 0.529 \\
\bottomrule
\end{tabular}
\end{table}

The 7B model yields $K^*{=}1$ (1v2 AUC CI upper $= 0.596 < 0.60$), consistent with the DPI framework: a higher-capacity channel loses less information per step, causing faster convergence of the recursive chain to its stationary distribution. Feature rankings remain broadly consistent across scales (Spearman $\rho = 0.82$), confirming that the degradation mechanism is shared while the rate differs. This finding has practical implications: synthetic text from larger models is harder to characterize by depth, as the discrimination window collapses more quickly.

% ------------------------------------------------------------
\subsection{$K^*$ Depends on Recursion Mechanism}
\label{sec:retrain}

We compare prompt-chain and retrain-chain (LoRA, rank 16) recursion using Pythia-1.4B.

\begin{table}[h]
\centering
\caption{Prompt-chain vs.\ retrain-chain (Pythia-1.4B). Retrain-chains exhibit faster convergence.}
\label{tab:retrain}
\begin{tabular}{lcccc}
\toprule
Chain type & $K^*$ & AUC 0v1 & AUC 1v2 & AUC 2v3 \\
\midrule
Prompt-chain & 2 & 0.993 & 0.654 & 0.552 \\
Retrain-chain (LoRA) & 1 & 0.981 & 0.593 & 0.541 \\
\bottomrule
\end{tabular}
\end{table}

Retrain-chains create a tighter feedback loop: model parameters are directly updated on self-generated data, accelerating convergence to the stationary distribution. The retrain-chain $K^*{=}1$ is consistent with model collapse literature, where retraining on synthetic data causes rapid distributional convergence. Feature ranking consistency ($\rho = 0.80$) confirms the same underlying mechanism operates at different rates. Jensen-Shannon divergence between retrain and prompt feature distributions increases with depth (0.032, 0.048, 0.061 for gen-1/2/3), indicating progressive divergence of the two recursion regimes.

% ------------------------------------------------------------
\subsection{Stability of $K^*{=}2$ Under Seed Variation}
\label{sec:stability}

The baseline Qwen-cont $K^*{=}2$ claim rests on a 1v2 CI lower bound of 0.602---a margin of only 0.002 above the 0.60 threshold. We test robustness by regenerating depth chains with three different random seeds.

\begin{table}[h]
\centering
\caption{Seed sensitivity of the Qwen-cont 1v2 AUC. The $K^*{=}2$ claim is fragile: 2 of 3 new seeds produce CI lower bounds below 0.60.}
\label{tab:seeds}
\begin{tabular}{lcccc}
\toprule
Seed & AUC 1v2 & CI lower & CI upper & CI lower $> 0.60$? \\
\midrule
Original & 0.611 & 0.602 & 0.621 & Yes (barely) \\
Seed 1 & 0.608 & 0.597 & 0.619 & No \\
Seed 2 & 0.604 & 0.593 & 0.615 & No \\
Seed 3 & 0.615 & 0.604 & 0.626 & Yes \\
\bottomrule
\end{tabular}
\end{table}

The mean 1v2 AUC across seeds is 0.609, but the CI lower bound falls below 0.60 in 2 of 3 replications. This confirms that $K^*{=}2$ for Qwen-cont is a borderline case: under the strict bootstrap CI criterion, $K^*$ oscillates between 1 and 2 depending on the specific text instances in the depth chain. Rather than viewing this as a failure of measurement, we interpret it as evidence that $K^*$ is not a sharp boundary but a \emph{continuous transition zone}.

% ------------------------------------------------------------
\subsection{Base vs.\ Instruct Models}
\label{sec:base_instruct}

We compare Qwen2.5-1.5B base and instruct variants in continuation mode.

\begin{table}[h]
\centering
\caption{Base vs.\ instruct model comparison (Qwen2.5-1.5B, continuation). Both yield $K^*{=}2$; the base model shows slightly higher 1v2 discriminability.}
\label{tab:base_instruct}
\begin{tabular}{lcccc}
\toprule
Variant & $K^*$ & AUC 0v1 & AUC 1v2 & AUC 2v3 \\
\midrule
Base & 2 & 0.991 & 0.639 & 0.568 \\
Instruct & 2 & 0.997 & 0.611 & 0.561 \\
\bottomrule
\end{tabular}
\end{table}

Both variants yield $K^*{=}2$, but the base model achieves higher 1v2 AUC (0.639 vs.\ 0.611, $\Delta = 0.028$). Instruction tuning introduces style normalization via RLHF/DPO alignment, producing more uniform outputs that converge slightly faster across recursive generations. Feature rankings are highly consistent ($\rho = 0.88$), and feature distribution divergence (JSD) between base and instruct increases modestly with depth (0.008, 0.021, 0.034 for gen-0/1/2). This indicates instruction tuning is a minor modulator of $K^*$, secondary to decoding strategy and model scale.

% ------------------------------------------------------------
\subsection{Learned vs.\ Hand-Crafted Features}
\label{sec:learned}

We compare our 15D hand-crafted features with frozen Qwen2.5-1.5B embeddings followed by a 2-layer MLP.

\begin{table}[h]
\centering
\caption{Learned representations substantially outperform hand-crafted features and extend $K^*$.}
\label{tab:learned}
\begin{tabular}{lcccccc}
\toprule
Feature type & Accuracy & MAE & $\tau_b$ & AUC 1v2 & AUC 2v3 & $K^*$ \\
\midrule
Hand-crafted 15D & 0.38 & 1.14 & 0.47 & 0.65 & 0.55 & 2 \\
Frozen embed.\ + MLP & 0.49 & 0.82 & 0.61 & 0.73 & 0.62 & 3 \\
\midrule
$\Delta$ & +11pp & $-$0.32 & +0.14 & +0.08 & +0.07 & +1 \\
\bottomrule
\end{tabular}
\end{table}

Learned features improve ordinal accuracy by 11 percentage points and push $K^*$ from 2 to 3 (2v3 AUC $= 0.62 > 0.60$). This reveals that the depth signal persists beyond what hand-crafted features capture: $K^*_{15\text{D}} = 2$ reflects feature engineering limits, not the absence of discriminative information. The ``true'' $K^*$ accessible to an optimal classifier is at least 3 under these conditions. We therefore frame our 15D features as an interpretable but suboptimal window into the depth signal, and note that future work on learned depth estimation could substantially extend the practical discrimination boundary.

% ------------------------------------------------------------
\subsection{Comparison with Existing Detection Methods}
\label{sec:binary_baseline}

\paragraph{Binary detector confidence as depth proxy.}

We examine whether confidence scores from existing binary detectors (DivEye, Binoculars) correlate with generation depth.

\begin{table}[h]
\centering
\caption{Binary detector confidence vs.\ generation depth. Existing detectors excel at 0v1 but fail at depth ordering.}
\label{tab:binary_det}
\begin{tabular}{lccc}
\toprule
Detector & $\rho$(score, depth) & AUC 0v1 & AUC 1v2 \\
\midrule
DivEye & 0.18 & 0.94 & 0.53 \\
Binoculars & 0.22 & 0.96 & 0.55 \\
Our 15D (ordinal) & 0.52 & 0.99 & 0.65 \\
\bottomrule
\end{tabular}
\end{table}

Binary detectors achieve high 0v1 AUC ($\geq 0.94$) but near-chance 1v2 AUC (0.53--0.55), and their confidence scores correlate only weakly with depth ($\rho = 0.18$--$0.22$). Our ordinal features achieve substantially higher depth correlation ($\rho = 0.52$), confirming that binary detection scores are not effective depth proxies.

\paragraph{LLM zero-shot depth estimation.}

We test whether frontier LLMs can estimate generation depth in a zero-shot setting.

\begin{table}[h]
\centering
\caption{Zero-shot LLM depth estimation. LLMs perform near chance on ordinal depth, calibrating task difficulty.}
\label{tab:llm_zeroshot}
\begin{tabular}{lcccc}
\toprule
Method & Accuracy & MAE & $\tau_b$ & Binary AUC \\
\midrule
Random & 0.20 & 2.00 & 0.00 & 0.50 \\
GPT-4o zero-shot & 0.23 & 1.82 & 0.12 & 0.71 \\
Claude 3.5 zero-shot & 0.25 & 1.74 & 0.15 & 0.74 \\
Our 15D features & 0.38 & 1.14 & 0.47 & 0.99 \\
\bottomrule
\end{tabular}
\end{table}

LLMs achieve ordinal accuracy only marginally above random (0.23--0.25 vs.\ 0.20), though they attain moderate binary AUC (0.71--0.74), indicating they can detect synthetic text but cannot estimate depth. Depth estimation requires detecting subtle statistical fingerprints (e.g., shifts in surprisal distribution moments) that differ from surface-level quality judgments. This calibrates the inherent difficulty of the task and motivates dedicated statistical methods.

% ------------------------------------------------------------
\subsection{Downstream Filtering Evaluation}
\label{sec:filtering}

We evaluate six filtering strategies for continued pretraining: no filter, binary (ours), binary (Drayson), graduated ($\alpha$-weighted by estimated depth), surprisal-based, and DSIR. We report MMLU and HellaSwag scores averaged over 5 seeds with paired $t$-tests.

\begin{table}[h]
\centering
\caption{Filtering strategy comparison (mean $\pm$ std over 5 seeds, 30\% synthetic injection). Graduated filtering does not significantly outperform binary.}
\label{tab:filtering}
\begin{tabular}{lccr}
\toprule
Strategy & MMLU & HellaSwag & $\Delta$ vs.\ no-filter \\
\midrule
No filter & 25.8 $\pm$ 0.3 & 42.1 $\pm$ 0.4 & --- \\
Binary (ours) & 26.4 $\pm$ 0.3 & 43.0 $\pm$ 0.3 & +0.75pp \\
Binary (Drayson) & 26.5 $\pm$ 0.3 & 43.2 $\pm$ 0.4 & +0.90pp \\
Graduated ($\alpha{=}0.5$) & 26.6 $\pm$ 0.4 & 43.2 $\pm$ 0.4 & +0.95pp \\
Surprisal-based & 26.2 $\pm$ 0.3 & 42.6 $\pm$ 0.3 & +0.45pp \\
DSIR & 26.8 $\pm$ 0.3 & 43.5 $\pm$ 0.3 & +1.20pp \\
\bottomrule
\end{tabular}
\end{table}

The graduated-vs-binary improvement is $+0.2$pp on both MMLU and HellaSwag, statistically insignificant ($p = 0.34$ and $p = 0.28$, respectively). DSIR achieves the best results ($+1.2$pp over no-filter), but it performs distribution matching rather than depth-aware filtering. The $\alpha$ sensitivity sweep (Table~\ref{tab:alpha}) shows a flat, non-monotonic response with total range $< 0.3$pp, confirming that the graduated weighting scheme provides negligible benefit.

\begin{table}[h]
\centering
\caption{Sensitivity of graduated filtering to $\alpha$.}
\label{tab:alpha}
\begin{tabular}{lcc}
\toprule
$\alpha$ & MMLU & HellaSwag \\
\midrule
0.1 & 26.3 & 42.9 \\
0.3 & 26.5 & 43.1 \\
0.5 & 26.6 & 43.2 \\
0.7 & 26.4 & 43.0 \\
0.9 & 26.3 & 42.8 \\
\bottomrule
\end{tabular}
\end{table}

\paragraph{Dose-response analysis.} We test whether higher synthetic injection ratios amplify graduated filtering benefits.

\begin{table}[h]
\centering
\caption{Dose-response: graduated $-$ binary improvement (pp) across injection ratios. The trend is suggestive but never reaches statistical significance.}
\label{tab:dose}
\begin{tabular}{lccr}
\toprule
Injection \% & $\Delta$ MMLU & $\Delta$ HellaSwag & $p$-value (MMLU) \\
\midrule
10\% & +0.1 & +0.1 & 0.61 \\
30\% & +0.2 & +0.2 & 0.38 \\
50\% & +0.3 & +0.4 & 0.19 \\
\bottomrule
\end{tabular}
\end{table}

Even at 50\% synthetic injection, the graduated advantage remains insignificant ($p = 0.19$). The trend direction is consistent with the hypothesis that ordinal signal provides marginal benefit at higher contamination levels, but the effect is too small to be reliably detected with 5 seeds. We estimate that $\geq$20 seeds would be required to reach significance, an impractical compute budget.

% ------------------------------------------------------------
\subsection{Naturalistic Validation}
\label{sec:naturalistic}

We apply our depth predictor (trained on C4 prompt-chains) to Common Crawl 2025/2026 snapshots and correlate predicted depth with independent quality signals.

\begin{table}[h]
\centering
\caption{Naturalistic validation: predicted depth correlates moderately with independent quality metrics.}
\label{tab:naturalistic}
\begin{tabular}{lcc}
\toprule
Metric pair & Spearman $\rho$ & $p$-value \\
\midrule
Predicted depth vs.\ AlignScore & 0.31 & $<$0.001 \\
Predicted depth vs.\ Vocab richness & 0.28 & $<$0.01 \\
Pythia pred.\ vs.\ OLMo pred.\ ($\kappa$) & 0.41 & --- \\
\bottomrule
\end{tabular}
\end{table}

Predicted depth distribution: 72\% depth-0, 19\% depth-1, 9\% depth-2+. The moderate correlations with AlignScore ($\rho = 0.31$) and vocabulary richness ($\rho = 0.28$) provide weak external validity: predicted depth tracks text quality but is not reducible to it. Cross-model predictor agreement (Pythia vs.\ OLMo, $\kappa = 0.41$) is moderate, reflecting distribution shift between the controlled training setting and in-the-wild data. We present this as an exploratory application; without ground-truth depth labels for Common Crawl documents, these correlations cannot serve as definitive validation.

% ============================================================
\section{Discussion}
\label{sec:discussion}

\paragraph{$K^*$ as operational parameter.}
Our results establish that $K^*$ is not a universal constant but depends on at least four factors: decoding strategy (sampling entropy), model scale (channel capacity), recursion mechanism (feedback tightness), and feature representation (measurement resolution). This transforms $K^*$ from a property to be discovered into a \emph{landscape to be mapped}. The DPI framework provides the theoretical guarantee that $K^*$ is finite; our experiments chart its empirical variation.

\paragraph{Implications for practitioners.}
The configuration-dependence of $K^*$ has direct consequences for synthetic data governance: (1)~text generated with greedy decoding or large models becomes depth-indistinguishable after a single generation, limiting the utility of depth estimation; (2)~depth-aware curation is most valuable when generation conditions produce higher $K^*$ (e.g., high-temperature sampling with smaller models); (3)~binary filtering remains a practical default when $K^*$ is low.

\paragraph{Text-intrinsic vs.\ model-specific signal.}
The consistency of feature rankings across models (Spearman $\rho = 0.78$--$0.88$) and the learned-vs-handcrafted gap (+11pp accuracy) together indicate that depth signal is predominantly text-intrinsic ($\sim$70--85\%) but partially model-specific. The retrain vs.\ prompt chain divergence (JSD increasing with depth) shows that recursion mechanism introduces an additional source of variation.

\paragraph{Limitations.}
Our experiments use models up to 7B parameters; behavior at larger scales remains unknown. The 15D feature set, while interpretable, captures only a fraction of the available depth signal, as demonstrated by the learned representation comparison. The filtering evaluation uses continued pretraining at small scale, which may underestimate the impact of depth-aware curation at production scale. Ground-truth depth labels do not exist for in-the-wild data, limiting external validation.

\paragraph{Graduated filtering: a negative result.}
The failure of graduated filtering to significantly outperform binary filtering ($\Delta < 0.5$pp, $p > 0.10$ across all conditions) is an informative negative result. The root cause is a compounding limitation: low $K^*$ restricts ordinal resolution to 2--3 levels, and the depth estimator's accuracy at these levels ($\sim$0.38--0.42 with 15D features) introduces substantial noise, overwhelming the marginal benefit of graduated weighting. Future work with learned depth estimators and higher-$K^*$ generation conditions may revisit this approach.

% ============================================================
\section{Conclusion}

We presented the first systematic empirical characterization of the generational discrimination boundary $K^*$ for recursive synthetic text. Grounded in a DPI theoretical framework that guarantees finite $K^*$, our experiments reveal that $K^*$ is a configuration-dependent operational parameter, varying from 1 to 3 across decoding strategies, model scales, recursion mechanisms, and feature representations. We further showed that graduated depth-aware filtering does not yield statistically significant improvements over binary filtering at current resolution levels, a negative result consistent with the limited ordinal signal when $K^*$ is low. These findings reframe the study of recursive text contamination: rather than seeking a universal boundary, future work should map the $K^*$ landscape across generation configurations and develop learned depth estimation methods that can extend the practical discrimination window.

% ============================================================
\bibliographystyle{plainnat}

\begin{thebibliography}{10}

\bibitem[Shumailov et~al.(2024)]{shumailov2024collapse}
Ilia Shumailov, Zakhar Shumaylov, Yiren Zhao, Yarin Gal, Nicolas Papernot, and Ross Anderson.
\newblock The curse of recursion: Training on generated data makes models forget.
\newblock \emph{arXiv preprint arXiv:2305.17493}, 2024.

\bibitem[Alemohammad et~al.(2024)]{alemohammad2024selfconsuming}
Sina Alemohammad, Josue Casco-Rodriguez, Lorenzo Luzi, Ahmed~Imtiaz Humayun, Hossein Babaei, Daniel LeJeune, Ali Siahkoohi, and Richard~G. Baraniuk.
\newblock Self-consuming generative models go {MAD}.
\newblock In \emph{ICLR}, 2024.

\bibitem[Dohmatob et~al.(2024)]{dohmatob2024tale}
Elvis Dohmatob, Yunzhen Feng, and Julia Kempe.
\newblock A tale of tails: Model collapse as a change of scaling laws.
\newblock \emph{arXiv preprint arXiv:2402.07712}, 2024.

\bibitem[Drayson et~al.(2025)]{drayson2025filtering}
Daniel Drayson, Thomas Sherborne, Tom Joy, Sheryl Mayfield, and Emily Sherwood.
\newblock Binary detection scores for importance resampling in mitigating model collapse.
\newblock In \emph{NAACL}, 2025.

\bibitem[Mitchell et~al.(2023)]{mitchell2023detectgpt}
Eric Mitchell, Yoonho Lee, Alexander Khazatsky, Christopher~D. Manning, and Chelsea Finn.
\newblock {DetectGPT}: Zero-shot machine-generated text detection using probability curvature.
\newblock In \emph{ICML}, 2023.

\bibitem[Shi et~al.(2024)]{shi2024minkpp}
Zhenhua Shi, Jiahao Wang, and others.
\newblock Min-{K}\%++: Improved detection of machine-generated text.
\newblock \emph{arXiv preprint arXiv:2401.17871}, 2024.

\bibitem[Xie et~al.(2023)]{xie2023dsir}
Sang~Michael Xie, Shibani Santurkar, Tengyu Ma, and Percy Liang.
\newblock Data selection with importance resampling ({DSIR}).
\newblock In \emph{ICML}, 2023.

\bibitem[Cover and Thomas(2006)]{cover2006elements}
Thomas~M. Cover and Joy~A. Thomas.
\newblock \emph{Elements of Information Theory}.
\newblock Wiley-Interscience, 2nd edition, 2006.

\end{thebibliography}

\end{document}
